% =====================================================================
% P1 Appendix A.1 --- Theorem 1 (identification under regime-dependent
% observation).
%
% Source of truth: paper/theorems.md v0.4 frozen (2026-10-21).
% This file is a fragment: it is meant to be \input under
%   \appendix \section{Proofs}
% before A.2 and A.3.  The standalone driver main-a12.tex sets the
% subsection counter to 0, so this fragment becomes A.1.
%
% Requirements: amsmath, amssymb, amsthm, array, longtable, and
% hyperref (optional, but the fixed-tag mechanism below keeps
% destinations unique when it is loaded).  Equation tags (4.1)--(4.2)
% follow the frozen source; do not renumber without updating the main
% text.
% =====================================================================

\allowdisplaybreaks

\providecommand{\TV}{\mathrm{TV}}

\newtheorem*{aITthm}{Theorem 1}
\newtheorem*{aIlemA}{Lemma 1}
\newtheorem*{aIlemB}{Lemma 2}
\newtheorem*{aIlemC}{Lemma 3}
\newtheorem*{aIlemD}{Lemma 4}

% Fixed equation tags do not advance the equation counter, so hyperref
% would give every display the same destination.  A dedicated counter
% keeps the destinations unique without changing the printed (frozen)
% tag numbers.  The original \theHequation is restored at the end.
\providecommand{\theHequation}{}
\let\aItheHequation\theHequation
\newcounter{aIKanchor}
\renewcommand{\theHequation}{aI.\arabic{aIKanchor}}
\newcommand{\aIfixedtag}[1]{\stepcounter{aIKanchor}\tag{#1}}

\subsection{Identification under regime-dependent observation (Theorem~1)}\label{app:a1}

This appendix gives the population identification proof of Theorem~1 and the
consistency argument for the model-based plug-in.  It is \emph{not} a
finite-sample rate result and \emph{not} a nonparametric identification theorem.
The proof chain is Lemma~1 (anchor and short-range dynamics), Lemma~2 (freeze
bridge; needed only in the non-Markov branch), Lemma~3 (remaining POMDP
components), and Lemma~4 (value and plug-in).  In the experimental Markov
environment, Lemma~2 is not used.

\paragraph{Scope and notation.}
We use the notation of \S2.  Time is discrete, $t=0,\dots,H$; the latent state is
$X_t=(S_t,Z_t)$, with $S_t>0$ the latent reference price and $Z_t$ the bias
state, stale anchor, position, and cash.  The calendar regime is
$r_t\in\mathcal R=\{\mathrm{OPEN},\mathrm{EXTENDED},\mathrm{WEEKEND},
\mathrm{HOLIDAY},\mathrm{HALT}\}$, known at $t$ and exogenous under A5.  The
transition is $X_{t+1}\sim T_{r_t}(\cdot\mid X_t,a_t)$, the observation is
$P_t\sim h_{r_t}(\cdot\mid X_t)$, the reward is
\[
R_t=\rho(X_t,a_t,X_{t+1})
=\frac{\text{fees}_t+\Delta\mathrm{IL}^{\mathrm{true}}_t-\text{cost}_t}{V_0},
\]
and the initial law is $X_0\sim\nu$.  The observable history is
$H_t^o=(P_{0:t},r_{0:t},a_{0:t-1},\text{position/cash})$, the belief is
$b_t(\cdot)=\mathbb P(X_t\in\cdot\mid H_t^o)$, and the target policy
$\pi(a_t\mid H_t^o)$ is known.  For an evaluation set
$R_{\mathrm{eval}}\subseteq\mathcal R$ fixed by the calendar,
\[
V^\pi(R_{\mathrm{eval}})
=\mathbb E_{\theta,\pi}\!\left[
\sum_{t:\,r_t\in R_{\mathrm{eval}}}\gamma^t\,\rho(X_t,A_t,X_{t+1})
\right],
\qquad \theta=(T,h,\nu,\rho).
\]
The full-window value takes $R_{\mathrm{eval}}=\mathcal R$.

For $\Delta=1,2$, let
$\mathcal A_{\mathrm{OPEN},\Delta}=\{t:r_t=r_{t+\Delta}=\mathrm{OPEN},\ t,t+\Delta
\text{ in the same day}\}$ be the calendar-adjacent OPEN pairs.  OPEN
observations are not contiguous in calendar time (overnight, weekend, and
holiday gaps), so all OPEN time-series moments below use only
$\mathcal A_{\mathrm{OPEN},1}\cup\mathcal A_{\mathrm{OPEN},2}$; a gap of
$\Delta$ calendar hours inside a day contributes $\phi^\Delta$, and no moment
is taken across a closed regime.

For a freeze event $e$ with pre- and post-anchors $t_e^-,t_e^+$, length
$\tau_e=t_e^+-t_e^-$, and pre-freeze volatility $\sigma_e$, write
\[
Q^{\mathrm{end}}_{\tau,\sigma}(ds^+\mid s^-)
=\mathcal L(S_{t_e^+}\in ds^+\mid S_{t_e^-}=s^-,\tau_e=\tau,\sigma_e=\sigma)
\]
for the endpoint conditional law, and
$Q^{\mathrm{path}}_{\tau,\sigma}(\cdot\mid s^-)$ for the full path law.  These
are different objects: I5 identifies only the former, while the latter requires
I5$'$ or a Markov $T$.

The observation kernels are those of \S2.2: OPEN is
$P_t=S_t e^{\varepsilon_t}$ with i.i.d.\
$\varepsilon_t\sim N(0,\sigma_o^2)$; EXTENDED is $P_t=S_t e^{\eta_t}$ with a
zero-started block AR(1) bias; WEEKEND and HOLIDAY use
$P_t=S_{\mathrm{last\,open},t}e^{\eta_t}$; and HALT is $P_t=P_{t-1}$.  The bias
is reset to zero at every OPEN step and is not stationary inside a window.

\paragraph{Assumptions.}
The main assumptions I1$'$ (shared latent dynamics and cost structure), I2
(anchored OPEN observations), I3 (parametric observation family with exogenous
AR(1) bias innovations), and, only in part (c)(ii), I5 (cross-event variation)
and I5$'$ (bridge condition) are stated in \S3.1.  Table~\ref{tab:a1-assumptions}
records the auxiliary and technical conditions used below.

\begin{longtable}{@{}>{\raggedright\arraybackslash}p{0.07\linewidth}>{\raggedright\arraybackslash}p{0.89\linewidth}@{}}
\caption{Auxiliary and technical assumptions for Theorem~1.}\label{tab:a1-assumptions}\\
\hline
ID & Condition \\
\hline
\endfirsthead
\hline
ID & Condition \\
\hline
\endhead
\hline
\endfoot
A1 & \textbf{Bounded reward.} $|\rho|\le R_{\max}$ a.s.  Otherwise the
integrability of $V^\pi$ must be assumed separately; a heavy-tailed $\rho$ can
make $\mathbb E[\rho]$ undefined, which is a failure of the estimand, not of the
bound. \\
A2 & \textbf{Known policy.} The target policy $\pi$ is known; the logging policy
$\mu$ is known or estimable. \\
A3$'$ & \textbf{Model-based support.} For every $(x,a)$ reachable by $\pi$, the
transition $T(\cdot\mid x,a)$ is identified from data (requires
$\mu(a\mid x)>0$, or structural knowledge).  The experimental environment
satisfies this automatically by A7 and deterministic position transitions. \\
A4 & \textbf{TV kernel topology.} $\sup_{x,a}\TV(T_n(\cdot\mid x,a),
T(\cdot\mid x,a))\to0$, $\sup_x\TV(h_n(\cdot\mid x),h(\cdot\mid x))\to0$,
$\TV(\nu_n,\nu)\to0$, and $\rho_n\to\rho$ bounded uniformly; then $V^\pi$ is
continuous.  Pointwise weak convergence is not enough.  Within the parametric
family, $\hat\theta\to\theta$ plus regularity implies this TV convergence;
general nonparametric estimators need not attain sup-TV, so A4 must not be read
as a nonparametric consistency assumption. \\
A5 & \textbf{Calendar-exogenous regimes.} $r_t$ is determined by the calendar
and does not depend on $S_t$ or on the action.  A price-triggered halt violates
A5; it must be restricted to scheduled holidays/weekends or handled by a
sensitivity analysis. \\
A6 & \textbf{Initial distribution.} $\nu$ is known or identified from initial
observations.  This appendix treats A6 as a pure assumption: if the initial
window starts in WEEKEND/HOLIDAY, deconvolving
$y_0=\log S_{\mathrm{last\,open}}+\eta_1$ identifies the pre-window anchor law,
not $\nu=\mathcal L(X_0)$; obtaining $\nu$ requires propagating through the
pre-window gap using the identified $T$ and the I1$'$ extrapolation.  The
experimental $\nu$ is a point mass, so this is not a numerical issue, but the
theorem does not claim the propagation. \\
A7 & \textbf{No market impact.} Actions do not change the pool price, volume, or
other traders' behavior; the $S$-component of $T$ is independent of $a$. \\
T1 & \textbf{Anchored-parameter identifiability.} Given $\nu$ from A6,
$(\psi,\sigma_o)\mapsto\mathcal L(y_{0:n}\mid\mathrm{OPEN})$ is injective for
$n$ large, where $\psi$ parametrizes $T$.  In the linear-Gaussian AR(1) case the
known-start moment profile in Lemma~1 gives a constructive proof (requires
$|\phi|<1$ and $\phi\ne0$).  For the jump family used in the experiments, T1 is
an identification condition, not a proved result; the numerical check reported
in the experiments is a check, not a proof. \\
T2 & \textbf{Joint endpoint characteristic function.} The \emph{joint}
characteristic function $\varphi_e(\xi,\zeta)$ of the endpoint measurement
errors $(\varepsilon^-,\varepsilon^+)$ is nonzero for every $(\xi,\zeta)$.
Independence of the two endpoint errors plus Gaussian marginals is sufficient;
marginal Gaussianity alone is not.  Used only in branch (ii). \\
T3 & \textbf{Support transfer.} The evaluation anchor law is absolutely
continuous with respect to the training anchor law $\Pi$; $(\tau,\sigma)$ lies
in the sample support, and in the continuous case the support contains a
neighborhood of the target value.  Values outside the support are extrapolation,
not identification.  Used only in branch (ii). \\
T4 & \textbf{Reward observability.} Fees, volume, and cost are observable;
$\Delta\mathrm{IL}^{\mathrm{true}}$ is a deterministic function of the $S$ path
and the position; $\rho$ is a known function once the identified components are
fixed.  The law of volume and fees under $\pi$ is invariant by A7, not by T4;
T4 gives observability only. \\
\hline
\end{longtable}

\begin{aITthm}[Identification under regime-dependent observation]
Assume I1$'$, I2, I3, A1, A2, A4--A7, T1, T4, and A3$'$ whenever $T$ depends on
actions and $\pi$ selects actions with no logging support.  In part (c)(ii)
additionally assume I5, I5$'$, T2, T3.  Then:

\emph{(a) Anchor observation.} $h_{\mathrm{OPEN}}$ is identified: $\sigma_o$ and
the conditional law $\mathcal L(S_t\mid H_t^o)$ on OPEN histories are identified;
$P_t$ is a noisy measurement of $S_t$, and $S_t$ is not pointwise identified.

\emph{(b) Short-range dynamics.} The $S$ component and the EXTENDED bias
component of the one-step transition $T$ are identified from calendar-adjacent
OPEN and EXTENDED observations.  The remaining regime-specific bias components
are completed in (d); the full joint kernel is not claimed before (d).

\emph{(c) Freeze-period bridge.} Two branches:

\emph{(i) Markov branch (the experimental environment).} If $T$ is Markov on the
freeze horizon and its $S$-component is identified as in (b), then for every
freeze length $\tau$, the path law
$Q^{\mathrm{path}}_{\tau,\sigma}(\cdot\mid s^-)$ is the $\tau$-step composition
of the identified $S$-component of $T$ started at $s^-$.  No I5/I5$'$, T2, T3, or
support restriction is needed in this branch; $\sigma$ enters only if the
identified $S$ component explicitly depends on it (it does not in the
experimental family).

\emph{(ii) Non-Markov branch.} If $T$ is not Markov on the freeze horizon, then
under I5 (together with T2 and T3) the endpoint law
$Q^{\mathrm{end}}_{\tau,\sigma}(\cdot\mid s^-)$ is identified for $\Pi$-a.e.\
$s^-$, with errors-in-variables and regime-specific endpoint noise
$h_{r(t_e^-)}$, $h_{r(t_e^+)}$.  The full path law $Q^{\mathrm{path}}$ is
identified only under the additional I5$'$ (parametric bridge, or direct path
identification as an assumption); endpoint identification alone is not enough
for path-dependent rewards.

\emph{(d) POMDP components.} The remaining regime-specific bias components of
$T$, the observation kernels $h_r$ for all $r\in\mathcal R$, the initial
distribution $\nu$, and the reward function $\rho$ are identified (for $\rho$,
using T4 and A7); consequently the belief process
$b_t(\cdot)=\mathbb P(X_t\in\cdot\mid H_t^o)$ is identified for every observable
history.

\emph{(e) Value and plug-in.} For any known target policy $\pi$ satisfying A2
and A3$'$, the value $V^\pi(R_{\mathrm{eval}})$ is identified.  Moreover, if
$\hat\theta\to\theta$ in the TV kernel topology of A4, then the model-based
plug-in $\hat V^\pi=V^\pi(\hat\theta)$ is consistent:
$\hat V^\pi\to V^\pi(\theta)$ in probability (or a.s., according to the mode of
$\hat\theta\to\theta$).
\end{aITthm}

\paragraph{Proof chain.}
Lemma~1 gives (a)--(b).  Lemma~2 gives (c)(ii); in branch (i), (c) follows from
Lemma~1 and the Markov composition.  Lemma~3 gives (d).  Lemma~4 gives (e).

\paragraph{What the theorem does not say.}
\begin{itemize}
\item It does not identify $S_t$ pointwise: each time has one noisy observation;
what is identified is the measurement model and the conditional law.
\item It is not a nonparametric result: I3 and T1 are premises.  A nonparametric
version needs stronger deconvolution and support conditions and is not claimed.
\item It does not identify the full non-Markov freeze path: branch (ii)
identifies the endpoint law; the path law requires I5$'$, whose ``direct path
identification'' option is an assumption, not a conclusion.
\item It is not a finite-sample guarantee: Theorem~1 is population-level;
finite-sample coverage is Theorem~2.
\item It does not claim necessity of I5: I5/I5$'$ appear only in (c)(ii); the
experimental Markov environment is branch (i), where I5 is redundant.
Theorem~4 gives a within-class non-identification example, not a general
necessity claim; the stress variant probes the role of cross-event variation
separately.
\item It is not a real-data conclusion: the empirical work is synthetic and the
regimes are calendar-exogenous; a price-triggered halt violates A5.
\item It does not make pure IPS/DR available: $\Delta\mathrm{IL}^{\mathrm{true}}$
is latent, so model-based evaluation is required by the estimand.
\item It does not claim that errors-in-variables correction is numerically
necessary in the current environment (reliability is approximately one); it
identifies the measurement-error model.
\end{itemize}

\begin{aIlemA}[Anchor observation and short-range transition]
Assume I1$'$, I2, I3, A6, and T1.  Then:
\begin{enumerate}
\item[(i)] $h_{\mathrm{OPEN}}$ (i.e.\ $\sigma_o$) and
$\mathcal L(S_t\mid H_t^o)$ on OPEN histories are identified (the
cross-regime belief is completed in Lemma~3);
\item[(ii)] the $S$ component of the one-step transition $T$, and the EXTENDED
bias component, are identified from calendar-adjacent OPEN and EXTENDED
observations.  The remaining regime-specific bias components are identified in
Lemma~3.
\end{enumerate}
\end{aIlemA}

\begin{proof}
\begin{enumerate}
\item On OPEN, $y_t:=\log P_t=\log S_t+\varepsilon_t$ with
$\varepsilon_t\sim N(0,\sigma_o^2)$ i.i.d.  Under I1$'$, $\log S_t$ is a
parametric Markov process with transition $T_\psi$ and known initial law $\nu$
(A6; the experiment has $\nu=\delta_{\mu_{\log}}$, so the process is not
stationary on the window).  Hence $\{y_t\}$ on OPEN is a hidden Markov model
with emission $N(0,\sigma_o^2)$, transition $T_\psi$, and known start.
\item T1 says $(\psi,\sigma_o)\mapsto\mathcal L(y_{0:n})$ is injective for $n$
large.  Therefore the observed OPEN law identifies $(\psi,\sigma_o)$.
\item \emph{Constructive linear-Gaussian case (no jumps), known start.}  Let
$X_t:=\log S_t-\mu_{\log}$ and
$X_{t+1}=\phi X_t+\sigma\xi_t$ with $\phi=1-\kappa$, and let $X_0=x_0$ be given
by $\nu$; the experiment has the point-mass start $x_0=0$
($\log S_0=\mu_{\log}$).  Hence $\{X_t\}$ is not stationary on the window; the
stationary ARMA(1,1) formulas are only the $t\to\infty$ limit.  Put
$V_t:=\mathrm{Var}(X_t)$, known from A6; in the experiment
\[
V_t=\frac{\sigma_{\mathrm{eff}}^2(1-\phi^{2t})}{1-\phi^2},
\qquad V_0=0.
\]
Then
\[
\mathrm{Cov}(y_t,y_{t+k})=\phi^k V_t\quad(k\ge1),
\qquad
\mathrm{Var}(y_t)=V_t+\sigma_o^2.
\]
\emph{Calendar caveat.}  OPEN is not contiguous in calendar time.  Use only
pairs in $\mathcal A_{\mathrm{OPEN},1}\cup\mathcal A_{\mathrm{OPEN},2}$: within a
day, $\Delta=1$ gives 51 pairs per path (8 complete OPEN days with 6 pairs each,
plus 3 pairs on day 7, which HALT splits into $09{:}00$--$10{:}00$ and
$13{:}00$--$15{:}00$), and $\Delta=2$ gives 41 pairs per path ($8\times5$, plus
day 7's single $13{:}00$--$15{:}00$); no closed-regime gap is used.  For a pair
separated by $\Delta$ calendar hours inside a day, replace $\phi$ by
$\phi^\Delta$.

\emph{Pair-aligned moments at the same start $t$} (any $t$ with
$t,t+1,t+2\in\mathrm{OPEN}$ and $V_t>0$; the experiment's first OPEN is $t=9$,
and under the point-mass start any $t\ge1$ works) identify the parameters
exactly, without stationarity:
\[
\phi=\frac{\mathrm{Cov}(y_t,y_{t+2})}{\mathrm{Cov}(y_t,y_{t+1})},
\qquad
\sigma_o^2=\mathrm{Var}(y_t)-\frac{\mathrm{Cov}(y_t,y_{t+1})}{\phi},
\]
\[
\sigma_{\mathrm{eff}}^2
=\mathrm{Var}(y_{t+1})-\sigma_o^2
-\phi^2\bigl(\mathrm{Var}(y_t)-\sigma_o^2\bigr).
\]
The $t$-profile $V_t=\mathrm{Cov}(y_t,y_{t+1})/\phi$ is over-identifying.  The
conditions $|\phi|<1$ and $\phi\ne0$ remain necessary: if $\phi=0$, $\log S$ is
i.i.d.\ from $t=1$ and only $\sigma_{\mathrm{eff}}^2+\sigma_o^2$ is identified;
the theorem excludes this degenerate case (or treats it as a known-i.i.d.\
model).  Invertibility is not needed for moment identification.

\emph{Estimation caveat.}  Pooling variance levels across starts mixes the
$V_t$ profile, and the ratio $\phi$ is exact only when the $\Delta=1$ and
$\Delta=2$ pairs share the same start set; align by start time or include the
$(1-\phi^{2t})$ profile.
\item \emph{Jumps and random volatility.}  With the rare jump component
$J_t\sim N(0,(3\sigma)^2)$ with probability $0.005$, the second moments identify
only the total innovation variance
\[
\sigma_{\mathrm{eff}}^2=\sigma^2+0.005\cdot9\sigma^2=1.045\sigma^2;
\]
separating the jump and diffusion variances requires higher moments or full-law
injectivity, which is exactly T1.  If $\sigma$ is drawn per path (lognormal in
the experiment), the population formulas use
$\mathbb E[\sigma_{\mathrm{eff}}^2]$ in place of $\sigma_{\mathrm{eff}}^2$.
\item \emph{EXTENDED (zero-started blocks).}  $y_t=\log S_t+\eta_t$ with $\eta$
independent of $S$ and of $\varepsilon$ at the innovation level (I3).  The
experiment resets $\eta=0$ at every OPEN step; an EXTENDED block's recursion
therefore starts from $\eta_0=0$ at the preceding OPEN, and its first EXTENDED
step is $u=1$ with $\mathrm{Var}(\eta_1)=\sigma_{\mathrm{ext}}^2$.  Within a
block $\eta$ is a zero-started AR(1), not a stationary one.  Let $u=1,2,\dots$
count steps into the block; then
\[
\mathrm{Var}(\eta_u)
=\frac{\sigma_{\mathrm{ext}}^2(1-\phi_{\mathrm{ext}}^{2u})}
{1-\phi_{\mathrm{ext}}^2},
\qquad
\mathrm{Cov}(\eta_u,\eta_{u+k})
=\phi_{\mathrm{ext}}^k\,\mathrm{Var}(\eta_u).
\]
Use only calendar-adjacent pairs inside an EXTENDED block (weekday blocks run
16:00--next day 08:00, 17 hours; the pre-holiday Thursday block is truncated at
24:00, 8 hours; no block starts Friday evening and no OPEN lies inside a block).
Since $V_t$ and $\phi$ are identified from OPEN,
\[
\mathrm{Cov}(y_t,y_{t+1})=\phi V_t+\phi_{\mathrm{ext}}\mathrm{Var}(\eta_u),
\qquad
\mathrm{Var}(y_t)=V_t+\mathrm{Var}(\eta_u),
\]
which gives $\mathrm{Var}(\eta_u)$ and
$\mathrm{Cov}(\eta_u,\eta_{u+1})$ for each $u$.  On the same $u$,
\[
\phi_{\mathrm{ext}}
=\frac{\mathrm{Cov}(\eta_u,\eta_{u+1})}{\mathrm{Var}(\eta_u)},
\qquad
\sigma_{\mathrm{ext}}^2
=\frac{\mathrm{Var}(\eta_u)(1-\phi_{\mathrm{ext}}^2)}
{1-\phi_{\mathrm{ext}}^{2u}};
\]
the $u$-profile is over-identifying.  This uses independence of the innovations,
not merely zero correlation, so that the joint transition factorizes as
$T_S\otimes T_\eta$; I3 supplies it.  If $S$ and $\eta$ are correlated, the cross
term does not vanish and joint state-space identification is required (open).
The stationary formulas
$\gamma_\eta(0)=\sigma_{\mathrm{ext}}^2/(1-\phi_{\mathrm{ext}}^2)$ and
$\gamma_\eta(1)=\phi_{\mathrm{ext}}\sigma_{\mathrm{ext}}^2/
(1-\phi_{\mathrm{ext}}^2)$ are only the $u\to\infty$ limit and must not be
written as the block's moments.
\item \emph{Conditional law (OPEN histories).}  Given the identified
$(\psi,\sigma_o)$, the Kalman filter (linear-Gaussian case) or the Bayesian
filter (general HMM) gives $\mathcal L(S_t\mid H_t^o)$ on OPEN histories as a
deterministic functional of identified components; the cross-regime belief is
completed in Lemma~3.
\end{enumerate}

\paragraph{Gap.}
For nonparametric $T$, T1 must be replaced by a deconvolution condition on the
pair $(y_t,y_{t+1})$; the inverse problem is ill-posed.  The paper restricts
Theorem~1 to the parametric family and states the nonparametric version as an
open extension.  For the jump family used in experiments, T1 is assumed rather
than proved.  The constructive argument is a known-start moment proof: it uses
the identified initial variance $V_0$ (A6) and the $(1-\phi^{2t})$ profile, not
stationarity; pooled estimators that ignore the start-time profile are biased.
The full joint $T$ also needs the stale-bias component from Lemma~3; Lemma~1
only identifies the $S$ and EXTENDED components.
\end{proof}

\begin{aIlemB}[Endpoint and path laws]
Assume the endpoint observation models from Lemma~1 (with regime-specific
noise), T2, and the mutual independence of endpoint errors stated in I2/I3.
\begin{enumerate}
\item[(a)] \emph{Markov branch.}  If $T$ is Markov on the freeze horizon and its
$S$-component is identified as in Lemma~1, then for each freeze length $\tau$,
$Q^{\mathrm{path}}_{\tau,\sigma}(\cdot\mid s^-)$ is the $\tau$-step composition
of the identified $S$-component of $T$ started at $s^-$.  No I5/I5$'$, T2, T3, or
support condition is needed.
\item[(b)] \emph{Non-Markov branch.}  Under I5 (and T2, T3), the endpoint law
$Q^{\mathrm{end}}_{\tau,\sigma}(\cdot\mid s^-)$ is identified for $\Pi$-a.e.\
$s^-$.  The path law $Q^{\mathrm{path}}_{\tau,\sigma}$ is identified only under
I5$'$(a) (parametric bridge) or I5$'$(b) (direct path identification, an
assumption).
\end{enumerate}
\end{aIlemB}

\begin{proof}
\emph{(a)} Immediate: under the Markov property, the conditional law of
$(S_{t_e^-+1},\dots,S_{t_e^+})$ given $S_{t_e^-}=s^-$ is obtained by composing
the identified $S$-component of $T$ exactly $\tau$ times; $\tau$ only selects
the composition length ($\sigma$ enters only if the identified $S$ dynamics
depends on it).  Endpoint deconvolution, T2, and the anchor-support condition
are not needed in this branch.

\emph{(b)} The population proof uses the joint characteristic function of the
endpoint errors.  Let
$Y_e^-=\log P_{t_e^-}$, $Y_e^+=\log P_{t_e^+}$, and let
$\varphi_e(\xi,\zeta)=\mathbb E[e^{i\xi\varepsilon^-+i\zeta\varepsilon^+}]$ be
the joint measurement-error characteristic function.  T2 and I2/I3 give
$\varphi_e\ne0$ everywhere and independence from $(s^-,s^+)$, so
\[
\Phi_Y(\xi,\zeta)
=\mathbb E[e^{i\xi\log s^-+i\zeta\log s^+}]\,\varphi_e(\xi,\zeta).
\]
Hence
\[
\Phi_s(\xi,\zeta):=\mathbb E[e^{i\xi\log s^-+i\zeta\log s^+}]
=\Phi_Y(\xi,\zeta)/\varphi_e(\xi,\zeta)
\]
is determined by the observed law.  By uniqueness of characteristic functions,
the joint law of $(\log s^-,\log s^+)$ is identified; by disintegration,
$Q^{\mathrm{end}}_{\tau,\sigma}(\cdot\mid s^-)$ is identified for $\Pi$-a.e.\
$s^-$.  Endpoint identification uses I5 (with T2 and T3); for the path law one
must add I5$'$(a) or I5$'$(b), since I5 alone does not identify
$Q^{\mathrm{path}}$.

\paragraph{Caveats.}
(i) Endpoint noise is regime-specific: $\tau=81$ has an EXTENDED pre-anchor with
$\sigma_\eta\approx0.0033$, not $\sigma_o=0.0005$.  (ii) Identification is
$\Pi$-a.e.\ and requires T3 for evaluation anchors, including a neighborhood in
$(\tau,\sigma)$ when the conditioning value is continuous.  (iii) $\sigma_e$ is
estimated, not known, and covariate measurement error in $\sigma_e$ is not
covered by the population argument.  (iv) The Gaussian deconvolution is
ill-posed in finite samples.  (v) T2 is a joint-CF condition, not a marginal one.
(vi) In the experimental Markov environment only branch (a) is used.
\end{proof}

\begin{aIlemC}[POMDP components]
Assume Lemma~1 (Lemma~2 is needed only for the non-Markov branch), I3 (exogenous
bias innovations with known block starts), A5, A6, T4, and A7.  Then the
remaining regime-specific bias components of $T$, the observation kernels $h_r$
for all $r$, $\nu$, and $\rho$ are identified; hence $b_t$ is identified.
\end{aIlemC}

\begin{proof}
\begin{enumerate}
\item \emph{OPEN.}  Lemma~1.
\item \emph{EXTENDED.}  $h_{\mathrm{EXT}}(P\mid S,\eta)=S\exp(\eta)$ with $\eta$
an exogenous AR(1) that is zero-started at each block onset (I3).  Use the same
block-$u$ profile as Lemma~1 step~5: after subtracting the identified $S$
moments $V_t$, the observed within-block moments give $\mathrm{Var}(\eta_u)$ and
$\mathrm{Cov}(\eta_u,\eta_{u+1})$, hence
\[
\phi_{\mathrm{ext}}
=\frac{\mathrm{Cov}(\eta_u,\eta_{u+1})}{\mathrm{Var}(\eta_u)},
\qquad
\sigma_{\mathrm{ext}}^2
=\frac{\mathrm{Var}(\eta_u)(1-\phi_{\mathrm{ext}}^2)}
{1-\phi_{\mathrm{ext}}^{2u}};
\]
the $u$-profile is over-identifying.  For the general parametric
$h_{\mathrm{EXT}}$ in I3, apply the same injectivity condition as T1 to the
observed law conditional on the identified $S$ law.
\item \emph{WEEKEND/HOLIDAY.}  $P=S_{\mathrm{last\,open}}\exp(\eta)$, so
$y_t=\log P_t=\log S_{\mathrm{last\,open}}+\eta_u$.  The anchor is constant in
$t$ within a path but random across paths, so the raw marginal moments contain
$\mathrm{Var}(\log S_{\mathrm{last\,open}})$; they are not the moments of
$\eta$.  The anchor cancels in first differences
$d_t:=y_{t+1}-y_t=\eta_{u+1}-\eta_u$.  With
$V_\eta(u):=\mathrm{Var}(\eta_u)$, $\phi:=\phi_{\mathrm{stale}}$,
$\sigma^2:=\sigma_{\mathrm{stale}}^2$ (and $|\phi|<1$; the unit-root case
$\phi=1$ is excluded),
\[
\mathrm{Var}(d_t)=(1-\phi)^2V_\eta(u)+\sigma^2,
\]
\[
\mathrm{Cov}(d_t,d_{t+1})
=(1-\phi)\bigl(\phi(1-\phi)V_\eta(u)-\sigma^2\bigr),
\qquad
\mathrm{Cov}(d_t,d_{t+2})=\phi\,\mathrm{Cov}(d_t,d_{t+1}).
\]
Hence, on the same $u$ (and provided
$\mathrm{Cov}(d_t,d_{t+1})\ne0$; otherwise the $u$-profile is used jointly),
\[
\phi=\frac{\mathrm{Cov}(d_t,d_{t+2})}{\mathrm{Cov}(d_t,d_{t+1})},
\qquad
\sigma^2=\phi\,\mathrm{Var}(d_t)-\mathrm{Cov}(d_t,d_{t+1}),
\]
\[
V_\eta(u)
=\frac{\mathrm{Cov}(d_t,d_{t+1})+(1-\phi)\mathrm{Var}(d_t)}
{(1-\phi)^2}.
\]
\emph{Estimation caveat.}  The $\sigma^2$ inverse is a difference of same-order
terms; with a large carried $V_0$ the early-$u$ triples are ill-conditioned.
Prefer larger-$u$ triples or a joint $u$-profile fit; population identification
is unaffected.

The block start is not stationary: a block that follows OPEN starts from
$\eta=0$ (the experiment's day-11 Friday-evening block), while a block that
follows EXTENDED carries the last EXTENDED state (the experiment's day-4
holiday block, whose initial variance follows from the identified
$(\phi_{\mathrm{ext}},\sigma_{\mathrm{ext}}^2)$ and the known block length).  In
general
\[
V_\eta(u)=\phi^{2u}V_0+\frac{\sigma^2(1-\phi^{2u})}{1-\phi^2}
\]
with $V_0$ known from the model/calendar (or identified from the $u$-profile);
the $u$-profile is over-identifying.  The cross-event decomposition (anchor law
plus bias) is a supplementary consistency check, not the identification route.
HOLIDAY uses the same stale channel with its own calendar anchor; in the
experiment the day-4 HOLIDAY, the following weekend, and Monday
00:00--08:00 form one continuous 81-hour stale block (same anchor and same
$(\phi,\sigma^2)$), so the carried start is inherited from the last EXTENDED
step.  The initial 9-hour block is zero-started by the model's initialization
and its anchor is $S_0$; its anchor law is the pre-window object discussed under
A6.
\item \emph{HALT.}  $P=P_{\mathrm{prev}}$ is a degenerate kernel; identified by
inspection.
\item \emph{$\nu$.}  A6, treated as a pure assumption (known or identified).  If
the initial regime is stale, deconvolving
$y_0=\log S_{\mathrm{last\,open}}+\eta_1$ (the first stale step; $\eta_0=0$)
identifies the pre-window anchor law, not $\nu=\mathcal L(X_0)$; getting $\nu$
needs propagation through the pre-window gap using the identified $T$ and the
I1$'$ extrapolation.  The experiment's $\nu$ is a point mass, so this is not a
numerical issue, but the theorem does not claim the propagation.
\item \emph{$\rho$.}  Fees, volume, and costs are observable (T4);
$\Delta\mathrm{IL}^{\mathrm{true}}$ is a deterministic function of the $S$ path
and position.  The $S$ process and the position dynamics (deterministic given
actions) are identified.  The law of the exogenous volume/fee process under
$\pi$ is invariant by A7 (no market impact); T4 alone gives observability, not
invariance.  Hence $\rho$ is identified.
\item \emph{Beliefs.}  The Bayesian filter with identified $\nu,T,h$ gives $b_t$
as a deterministic functional of the observable history.  For general kernels
this uses regular conditional probabilities and a domination condition; in the
Gaussian/parametric family used here it is automatic.  A degenerate HALT
observation performs no update.
\end{enumerate}

\paragraph{Gap.}
The moment steps assume the observation bias is independent of $S$ and actions
at the innovation level, with independent innovations across blocks (I3); the
carried block-start state is a known function of past innovations, not a fresh
independent draw, so the decomposition is applied block by block.  If the bias
is correlated with $S$, the cross term $\gamma_{S\eta}$ does not vanish and
joint identification requires a joint state-space argument; the synthetic
environment satisfies independence.  The nonparametric $h_r$ version is open.
A6 is an assumption, not a proved identification of $\nu$ for a stale initial
window.  The stale within-block moments are computed from first differences
(which cancel the random anchor) together with the zero-start/carried-start
$u$-profile, not from raw variances or stationary autocovariances; ignoring the
$u$-profile biases pooled estimates.
\end{proof}

\begin{aIlemD}[Value identification and plug-in consistency]
Assume Lemma~3, A1, A2, A4, A7 (and A3$'$ if $T$ depends on actions and $\pi$
uses unsupported actions).  Then:
\begin{enumerate}
\item $V^\pi(\theta)$ is uniquely determined by $\theta$ and the known $\pi$;
\item if $\hat\theta\to\theta$ in the TV kernel topology of A4, then
$\hat V^\pi=V^\pi(\hat\theta)\to V^\pi(\theta)$;
\item pure IPS/DR is not available; A3 is neither an identification condition
nor a substitute for the model.
\end{enumerate}
\end{aIlemD}

\begin{proof}
The proof has four steps.

\begin{enumerate}
\item \emph{Identification.}  $V^\pi(\theta)$ is a functional of $\theta$ and
$\pi$: the belief recursion gives $b_t$; $\pi(a\mid H_t^o)$ is known; hence the
joint law of $(X_t,A_t)$ is determined by $\nu,T,h,\pi$, and $\rho$ is given.
$\Delta\mathrm{IL}^{\mathrm{true}}$ is a deterministic function of the $S$ path
and position, so it does not need to be observed in the log.
\item \emph{Continuity.}  Write
$\mathcal L^t:=\mathcal L(X_0,P_0,A_0,\dots,X_t,P_t,A_t)$ (and similarly
$\mathcal L_n^t$).  With
\[
\begin{aligned}
\Delta_T(n)&=\sup_{x,a}\TV(T_n,T),\\
\Delta_h(n)&=\sup_x\TV(h_n,h),\\
\Delta_\nu(n)&=\TV(\nu_n,\nu),\\
\Delta_\rho(n)&=\sup_z|\rho_n-\rho|,
\end{aligned}
\]
induction gives
\begin{equation}\label{eq:a1-4.1}
\TV(\mathcal L_n^t,\mathcal L^t)
\le \Delta_\nu(n)+(t+1)\bigl(\Delta_T(n)+\Delta_h(n)\bigr).
\aIfixedtag{4.1}
\end{equation}
The $(t+1)$ counts $t$ transitions and $t+1$ emissions.  Since
$|\rho|\le R_{\max}$,
\[
|\mathbb E_n[R_t]-\mathbb E[R_t]|
\le 2R_{\max}\TV(\mathcal L_n^{t+1},\mathcal L^{t+1})+\Delta_\rho(n),
\]
because $R_t$ depends on $X_{t+1}$ and hence on $\mathcal L^{t+1}$; this is the
source of the $(t+2)$ factor below.  Summing over $t=0,\dots,H$ gives
\begin{equation}\label{eq:a1-4.2}
|V_n^\pi-V^\pi|
\le (H+1)\Delta_\rho(n)
+2R_{\max}\sum_{t=0}^{H}\bigl[\Delta_\nu(n)+(t+2)(\Delta_T(n)+\Delta_h(n))\bigr]
\to0.
\aIfixedtag{4.2}
\end{equation}
For the infinite-horizon $\gamma<1$ extension, truncate at $H_N$ with tail
$2R_{\max}\gamma^{H_N+1}/(1-\gamma)$, apply \eqref{eq:a1-4.2}, then let
$N\to\infty$.  The main text uses finite $H$ and does not need this truncation.
\item \emph{Consistency.}  $\hat\theta\to\theta$ in the A4 topology plus the
deterministic inequality \eqref{eq:a1-4.2} gives $\hat V^\pi\to V^\pi$.
\item \emph{Why not IPS/DR.}  The IPS estimand requires observed $R_i$; here
$R_t=(\text{fees}+\Delta\mathrm{IL}^{\mathrm{true}}-\text{cost})/V_0$ and
$\Delta\mathrm{IL}^{\mathrm{true}}$ is latent.  Replacing it by a model
prediction is already model-based (possibly with IPS as a hybrid), not pure OPE.
\end{enumerate}

\paragraph{Caveats.}
(i) Identification is for the ideal $\theta$; plug-in with $\hat S$ adds an
errors-in-variables error not covered by Lemma~4.  (ii) A4 is a strong TV
topology; in the parametric family it follows from $\hat\theta\to\theta$ under
regularity, but general nonparametric estimators do not attain sup-TV, so (e)
should not be read nonparametrically.  (iii) A1 fails under heavy tails, in
which case $V^\pi$ may not exist.  (iv) A3$'$ is needed if $T$ depends on
actions without logging support.
\end{proof}

\paragraph{Proof status.}
Lemma~1 (the known-start linear-Gaussian/no-jump case, with T1 assumed for the
jump family) and Lemma~2(b) have derivations in Appendix A.1; Lemma~3 (exogenous
bias innovations, known block starts) has a derivation here as well, while the
nonparametric $h_r$ version and the correlated-bias case are open, and A6 is
an assumption.  Lemma~4 has a complete
population/consistency proof; it gives no rate, A4's sup-TV is attainable in the
parametric family, and the errors-in-variables error from using $\hat S$ is not
covered.  The logic chain is closed under the stated assumption set: the Markov branch is
Lemma~1 $\to$ Lemma~3 $\to$ Lemma~4, and the non-Markov branch is
Lemma~1 $\to$ Lemma~2 $\to$ Lemma~3 $\to$ Lemma~4.  The open items above are
explicitly not covered by the theorem.

\let\theHequation\aItheHequation
